\Needspace{12\baselineskip}\SourceLeaf{315}

\begin{center}
[2.]
\end{center}

\begin{center}
\textsc{Gauss} an \textsc{Gerling}. Göttingen, 7. November 1830.
\end{center}

\indent\mbox{. . . . . . . . . .} Über die vier bemerkenswerten Punkte im Dreieck
habe ich selbst nie etwas veröffentlicht. Was ich jedoch 1808 darüber
\textsc{Schumacher} mitgeteilt habe, hat dieser im Anhang zu seiner
Übersetzung von \textsc{Carnot}, Géométrie de position, zweiter
Bd., der 1810 erschienen ist, drucken lassen. Ich bin versucht zu bezweifeln, dass
der Punkt, der sich nach dem von Ihnen angeführten Satz ergibt,
besonders elegante Beziehungen zu den übrigen aufweisen wird, und zwar
deshalb, weil jener Satz nur ein ganz spezieller Fall eines viel
allgemeineren ist.

Legt man nämlich an die Seiten eines gegebenen Dreiecks $ABC$
drei andere $AB\gamma$, $\alpha BC$, $AC\beta$ (entweder alle drei
außen oder alle drei innen am Dreieck
$ABC$), die jeweils einem zweiten gegebenen Dreieck $abc$ ähnlich
sind, und zwar so, dass sie auch ähnlich liegen und die Winkel bei
$\alpha$, $\beta$, $\gamma$ denen bei $a$, $b$, $c$ gleich sind, so
schneiden sich die Geraden $A\alpha$, $B\beta$, $C\gamma$
(nötigenfalls vorwärts oder rückwärts verlängert) in ein und demselben Punkt
$D$, und die Winkel bei $D$ ($ADB$, $BDC$, $CDA$) sind
entweder den Winkeln bei $\gamma$, $\alpha$, $\beta$ selbst oder ihren
Nebenwinkeln gleich und liegen ähnlich.

Offenbar ist dies die eleganteste Konstruktion der
Aufgabe von \textsc{Pothenot}. In der Praxis kann man sich aber damit begnügen,
nur einen Punkt, z.\,B. $\alpha$, zu konstruieren und den Messtisch nach
der Geraden $A\alpha$ zu orientieren. Wie ich von meinem Sohn höre, ist
dieses Verfahren unter dem Namen der \textsc{Schulz-Montanus}'schen Methode
bekannt. \mbox{. . . . . . .}
